\section{Conditional recourse: a small nonconvex core}\label{sec:rec}

Many structured problems contain a few complicating continuous variables
and many variables that are easy to optimize once the complicating ones are
fixed. Two-stage models with a few first-stage decisions, design problems
with a few shape or price parameters, and decomposition methods that branch
only on linking variables have this form. We call the complicating variables
the \emph{core} $v\in[0,1]^k$ and the others the \emph{residual} variables
$z$. Of the three structural routes in this paper, a small nonconvex factor,
sparse bags, and a small core with tractable conditional recourse, this
section and \cref{sec:int} develop the third. Its parameter is the core
size $k$ together with the ratio $L/\sigma$ of core curvature to noise. The
number of residual variables, their coupling, and their curvature enter only
through the input length.

The mechanism has four parts.
\begin{enumerate}[label=(\roman*)]
\item Minimizing over a residual set that does not depend on the core
preserves upper coordinate curvature. The conditional value
$V(v)=\min_zF_\gamma(v,z)$ therefore obeys the corrected-corner inequality
of \cref{thm:count:local} on the core cube, and no residual coordinate is
rounded.
\item Independent linear noise on the core coefficients bounds the expected
number of near-optimal core grid points by a product of $k$ one-dimensional
factors. Residual noise, when present, is conditioned on.
\item Residual coordinates are never refined, so closure needs a separate
global certificate. A few restricted recourse calls prove that every global
optimizer lies in a small product patch (this section) or uses one residual
integer label (\cref{sec:int}).
\item A base-selected finite law makes failure of these certificates rare
enough to pay for an exact fallback on the same draw
(\cref{sec:model:algorithms}).
\end{enumerate}
\Cref{tab:rec:summary} lists the results of this section. Throughout,
\begin{equation}\label{eq:rec:Q}
 Q_{\rm ex}=\Bigl[3+\frac{(1+k/2)L}{2\sigma}\Bigr]^k,
 \qquad
 Q_{\rm ap}=\Bigl[3+\frac{(1+k)L}{2\sigma}\Bigr]^k .
\end{equation}
The first factor applies when conditional values are computed exactly, the
second when they are computed with certified error. Both are numerical:
they depend on the magnitude of $L/\sigma$, not on its encoding length.

\begin{table}[t]
\centering
\caption{Continuous conditional recourse. Here $b=\log_2M$ is the number of
sampling bits per perturbed coordinate, ``implicit'' is the patch output of
\cref{def:model:outputs}\ref{it:model:implicit}, and every result returns an
algebraic output from the same-draw fallback on rare draws.}
\label{tab:rec:summary}
\small
\begin{tabularx}{\textwidth}{@{}lXlll@{}}
\toprule
Result & Noise; residual premise and oracle & $b$ & Expected work & Ordinary output\\
\midrule
\cref{thm:rec:qp} & ambient; quadratic, exact recourse on every residual
subbox & $\poly(I)$ & $8^kQ_{\rm ex}\poly(I)$ & rational\\
\cref{cor:rec:forest} & ambient; quadratic, residual interaction graph a
forest & $\poly(I)$ & $8^kQ_{\rm ex}\poly(I)$ & rational\\
\cref{thm:rec:poly} & ambient; certified values on every residual subbox,
e.g.\ residual convexity & $\poly_d(I)$ & $8^kQ_{\rm ap}\poly_d(I)$ & implicit\\
\cref{thm:rec:core-only} & core-only; residual Hessian $\succeq\mu I$
& $\poly_d(I)$ & $8^kQ_{\rm ap}\poly_d(I)$ & implicit\\
\bottomrule
\end{tabularx}
\end{table}

\subsection{The conditional value function}\label{sec:rec:value}

\begin{definition}[Core--recourse instance]\label{def:rec:instance}
A \emph{core--recourse instance} consists of integers $k,n_R\ge0$ with
$n=k+n_R\ge1$, a nonempty compact \emph{residual set}
$Y\subseteq\R^{n_R}$, an explicit rational polynomial $F_0(v,z)$ of total
degree at most $d$ in the core variables $v\in\R^k$ and residual variables
$z\in\R^{n_R}$, a rational $L>0$ with
\begin{equation}\label{eq:rec:L}
 \frac{\partial^2F_0}{\partial v_i^2}(v,z)\le L
 \qquad\bigl(i=1,\ldots,k,\ (v,z)\in[0,1]^k\times\conv Y\bigr),
\end{equation}
and a rational noise half-width $\sigma>0$. The feasible set is
$X=[0,1]^k\times Y$. We write $C$ and $R$ for the core and residual index
sets. In this section $Y=[0,1]^{n_R}$; \cref{sec:int} uses integer sets.
\end{definition}

The essential restriction is that the residual feasible set does not depend
on the core: the core enters only the objective. Core coordinates with other
rational ranges can be mapped affinely onto $[0,1]$; this rescales $L$ and
the physical noise, and the rescaled data are part of the input. A valid $L$
can always be computed from monomial coefficients, but the numerical factors
below use the supplied value, which may be much smaller; a sharper value is
a certified premise whose proof is charged to $I$. For the sampled
objective $F_\gamma(v,z)=F_0(v,z)+\gamma_C^{\mathsf T}v+\gamma_R^{\mathsf T}z$
(with $\gamma_R=0$ under core-only noise) define, for a nonempty closed
$Y'\subseteq Y$,
\begin{equation}\label{eq:rec:V}
 V_{Y'}(v)=\min_{z\in Y'}F_\gamma(v,z),\qquad V=V_Y,\qquad
 f^*=\min_XF_\gamma=\min_{[0,1]^k}V.
\end{equation}

\begin{lemma}[Conditional values]\label{lem:rec:value}
Let $G\ge\sup\{\sum_{i\in C}\abs{\partial_{v_i}F_\gamma(v,z)}:
v\in[0,1]^k,\ z\in\conv Y\}$.
\begin{enumerate}[label=(\alph*)]
\item The function $V(v)-\gamma_C^{\mathsf T}v$ depends on $\gamma_R$ but not
on $\gamma_C$.
\item For every $i\in C$ and every $v$, the function
$t\mapsto V(v+te_i)-Lt^2/2$ is concave on the interval where
$v+te_i\in[0,1]^k$.
\item For every nonempty closed $Y'\subseteq Y$, the function $V_{Y'}$ is
$G$-Lipschitz with respect to the maximum norm on the core:
$\abs{V_{Y'}(v)-V_{Y'}(w)}\le G\norm{v-w}_\infty$.
\end{enumerate}
None of these statements uses convexity of the residual problem.
\end{lemma}

Part~(b) is the reason for fixed residual feasibility. It makes $V$ satisfy
the curvature hypothesis of \cref{thm:count:local} in dimension $k$: for a
dyadic cell of side $h_j=2^{-j}$ in the core cube,
\begin{equation}\label{eq:rec:corner}
 \min_{v\in C}V(v)\ \ge\ \min_{v\text{ corner of }C}V(v)-E_j,
 \qquad E_j=\frac{kLh_j^2}8 .
\end{equation}
Only core coordinates contribute to $E_j$. When residual feasibility depends
on the core, \eqref{eq:rec:corner} fails, even for linear data.

\begin{example}[Core-dependent feasibility]\label{ex:rec:feasibility}
Let $k=n_R=1$, $F_0(v,z)=z$, and replace the fixed residual set by
$Y(v)=\{z\in[0,1]:z\ge v-\frac13,\ z\ge\frac13-v\}$. Then $\partial_{vv}F_0=0$
and $V(v)=\abs{v-\frac13}$. At every level $j$ the cell containing $\frac13$
has corners at distances $h_j/3$ and $2h_j/3$ from $\frac13$, because
$2^j\not\equiv0\pmod3$. Its minimum corner value is $h_j/3$, while
$\min_CV=0$. Thus \eqref{eq:rec:corner} fails with $L=0$, and for every
$L>0$ it fails at all levels with $h_j<8/(3kL)$. A search based on it would
report the value interval $[U_j-E_j,U_j]$, which excludes $f^*=0$.
\end{example}

The corrected-corner inequality also shows why conditional values must be
exact or certified, rather than computed by rounding the residual variables
to a grid. Rounding residual variables reintroduces a discretization error
whose variation over core states grows with the number of residual
variables, and normalizing grid messages does not remove it.

\begin{example}[Rounded residuals]\label{ex:rec:star}
Let $m=32$, $\varepsilon=\frac1{16}$, and on $[0,1]^{m+1}$ let
\[
 F_0(v,z)=v^2-\frac v2\sum_{i=1}^mz_i+\sum_{i=1}^mz_i^2
 +\Bigl(\frac m{16}-1-\varepsilon\Bigr)v .
\]
Its interaction graph is a star, every coordinate second derivative equals
$2$, and $V(v)=(1-\frac m{16})v^2+(\frac m{16}-1-\varepsilon)v$ is strictly
concave with $V(0)=0>V(1)=-\varepsilon$. The unique optimizer is
$(1,\frac14,\ldots,\frac14)$ with value $-\frac1{16}$. With $h=\frac12$, let
$V_h(v)$ minimize over residual points of the grid $\{0,\frac12,1\}^m$. Then
$V_h(0)=0$, $V_h(\frac12)=\frac{23}{32}$, and $V_h(1)=\frac{31}{16}$, and
$V_h(v)-V(v)=m\dist(v/4,\{0,\frac12,1\})^2$. The bag cell
$[\frac12,1]\times[0,\frac12]$ of $\{v,z_1\}$ contains the optimizer's bag
projection, its smallest grid min-marginal is $\frac{23}{32}$, the grid
incumbent is $0$, and the bag-local correction $\abs BLh^2/8=\frac18$
discards it. The global correction $(m+1)Lh^2/8=\frac{33}{16}$ is valid. With
exact conditional values $V$ in place of $V_h$, the corner correction is
$Lh^2/8$ for the single core coordinate, independently of $m$.
\end{example}

The proof (\cref{app:rec:examples}) also gives a uniformly positive definite
variant on a star with $m=d^2$ leaves, $d\ge8$, in which the variation of
$V_h-V$ between two adjacent core states is $(d-1)h^2$. Positive
definiteness and bounded coloring number therefore do not make a bag-local
budget valid for grid min-marginals. These examples concern the retention
rule; they are not hardness results.

\paragraph{The core search.}
A \emph{conditional oracle} receives a rational core point $v$. In
\emph{exact mode} it returns $V(v)$ and a residual point $z\in Y$ with
$F_\gamma(v,z)=V(v)$. In \emph{certified mode} it also receives a rational
accuracy $\eta>0$ and returns rationals $\ell\le V(v)\le u$ with
$u-\ell\le\eta$ and a residual point $z\in Y$ with $F_\gamma(v,z)=u$. For
uniform notation, exact mode returns $\ell=u=V(v)$.

The \emph{core search} processes levels $j=0,1,\ldots$. Level $0$ has the
single cell $[0,1]^k$. At level $j\ge1$ the search generates the $2^k$
children of every cell retained at level $j-1$; only these cells are
generated, and the full level-$j$ grid is never enumerated. It queries the
oracle once at every distinct corner of a generated cell, in certified mode
with $\eta=E_j$. Let $U_j$ be the least upper value $u$ returned at this
level, attained at a corner $c_j$ with completion $z_j$. A generated cell is
retained if and only if
\begin{equation}\label{eq:rec:retain}
 \min_{v\text{ corner of }C}\ell(v)-E_j\ \le\ U_j .
\end{equation}
The \emph{retained hull} $\mathcal H_j$ is the smallest box containing all
retained cells. For $k=0$ there is one empty core cell; each level makes one
query with accuracy $4^{-j}$ in certified mode, and $\mathcal H_j$ is a point.

\begin{proposition}[Core search]\label{prop:rec:search}
Let $\varepsilon_j=E_j$ in exact mode and $\varepsilon_j=2E_j$ in certified
mode. On every draw and for every admissible oracle answer:
\begin{enumerate}[label=(\alph*)]
\item for every global minimizer $(a,z^*)$ of $F_\gamma$ on $X$, a retained
cell contains $a$, so $a\in\mathcal H_j$;
\item $f^*\le U_j\le f^*+\varepsilon_j$, and $c_j$ is a corner of a retained
cell, so $c_j\in\mathcal H_j$;
\item every retained cell has a corner $v$ with $V(v)\le f^*+2\varepsilon_j$.
\end{enumerate}
Suppose the core coefficients $\gamma_C$ are independent with law
$U_{\sigma,M}$ and independent of $\gamma_R$, and $M\ge2^J$. Conditionally
on $\gamma_R$, for every $j\le J$ the expected number of points $v$ of the
full grid $(h_j\Z)^k\cap[0,1]^k$ with $V(v)\le f^*+2\varepsilon_j$ is at most
$Q_{\rm ex}$ in exact mode and $Q_{\rm ap}$ in certified mode. The expected
number of oracle calls through level $J$ is at most $2^k+8^kJQ$, with
$Q=Q_{\rm ex}$ or $Q=Q_{\rm ap}$ respectively.
\end{proposition}

The expectation statement uses a deterministic grid fixed before sampling.
Each retained cell is charged to a near-optimal corner, and near-optimality
of a fixed grid point is a condition on the true function $V$, so arbitrary
noise-dependent choices of the oracle within its error bounds do not affect
the count. This is the conditional form of \cref{thm:count:local}; the proof
in \cref{app:rec:search} is self-contained.

\subsection{Global containment and closure certificates}\label{sec:rec:certificates}

The core search alone does not finish the problem: it localizes the core but
says nothing about residual coordinates. The next two lemmas supply sound
certificates. Here $Y=\prod_{i\in R}[\ell_i,u_i]$ is a box.

Let $c$ be a point of a core box $\mathcal H$, let $z_c\in Y$, and let $r>0$.
The \emph{residual patch} is
\[
 P=\prod_{i\in R}\bigl[\max\{\ell_i,z_{c,i}-r\},\ \min\{u_i,z_{c,i}+r\}\bigr].
\]
For each $i\in R$ with $\ell_i<z_{c,i}-r$ form the \emph{lower slab}
$Y_i^-=\{z\in Y:z_i\le z_{c,i}-r\}$, and for each $i\in R$ with
$z_{c,i}+r<u_i$ the \emph{upper slab} $Y_i^+=\{z\in Y:z_i\ge z_{c,i}+r\}$.
No slab is formed on a side where the patch reaches an original bound; a
singleton slab at an original bound could contain an optimizer and destroy
the certificate.

\begin{lemma}[Excluded slabs]\label{lem:rec:exclusion}
The slabs cover $Y\setminus P$. Let $\lambda_{\rm out}$ be at most the least
conditional value $V_{Y_i^\pm}(c)$ over all slabs ($\lambda_{\rm out}=+\infty$
if there are none), let $U\ge V(c)$, and let $w_\infty(\mathcal H)$ be the
largest coordinate width of $\mathcal H$. If
\begin{equation}\label{eq:rec:exclusion}
 \lambda_{\rm out}-U>2G\,w_\infty(\mathcal H),
\end{equation}
then for every $v\in\mathcal H$ every minimizer of $F_\gamma(v,\cdot)$ on $Y$
lies in $P$. If moreover $\mathcal H$ contains the core of every global
minimizer, then every global minimizer of $F_\gamma$ on $X$ lies in
$\mathcal H\times P$.
\end{lemma}

The certificate needs a certified \emph{lower} bound on every slab. A
feasible value on a slab, which is all an approximate solver returns by
itself, does not suffice: a better competitor could be hidden in the slab.
It also needs restricted recourse on slabs, so the conditional oracle must
accept arbitrary residual subboxes. This is the \emph{box-stability}
requirement of the following sections.

Once global containment is certified, two sound local tests finish the
problem. Let $Q\subseteq[0,1]^n$ be a box containing every global
minimizer, with midpoint $m$ and maximum-norm radius $r_Q$, and put
\begin{equation}\label{eq:rec:K}
 K_2\ge\max\Bigl\{1,\max_i\sum_j\sup_{[0,1]^n}\abs{\partial_{ij}F_0}\Bigr\},
 \qquad
 K_3\ge\max\Bigl\{1,\max_i\sum_{j,l}\sup_{[0,1]^n}\abs{\partial_{ijl}F_0}\Bigr\},
\end{equation}
rational bounds computed from monomial coefficients.

\begin{lemma}[Sign fixing and convexity tests]\label{lem:rec:closure}
\leavevmode
\begin{enumerate}[label=(\alph*)]
\item If the original lower bound of coordinate $i$ lies in the $i$th
interval of $Q$ and $\partial_iF_\gamma>0$ on $Q$, every global minimizer has
$x_i$ at that bound; symmetrically for an upper bound and
$\partial_iF_\gamma<0$. The midpoint test $\pm\partial_iF_\gamma(m)>K_2r_Q$
is a sufficient condition for the sign; for quadratics the sign can be
decided exactly. Every global minimizer lies in the box $Q'$ obtained by
fixing all coordinates that pass.
\item Let $m'$ and $r_{Q'}$ be the midpoint and radius of $Q'$ and let
$g_0>0$. If the Hessian of $F_0$ in the unfixed coordinates satisfies
$\nabla^2F_0(m')-(K_3r_{Q'}+g_0)I\succeq0$, then $F_\gamma$ is
$g_0$-strongly convex on $Q'$, its unique minimizer on $Q'$ is the unique
global minimizer of $F_\gamma$ on $X$, and $(Q',g_0)$ is an implicit patch
output (\cref{def:model:outputs}).
\item If $F_0$ is quadratic and its Hessian in the unfixed coordinates is
positive semidefinite, then every exact minimizer of the convex quadratic
program $\min_{Q'}F_\gamma$ is a global minimizer of $F_\gamma$ on $X$.
\end{enumerate}
\end{lemma}

The tests may fix only \emph{original} bounds, and only bounds that lie in
the current interval. Fixing an artificial patch endpoint would be unsound.
A coordinate whose derivative vanishes at a contained optimizer never passes
a strict sign test. All tests in this section are exact rational
computations and are sound on every draw; the probabilistic analysis only
shows that they pass early.

\subsection{Exact recourse for quadratic objectives}\label{sec:rec:qp}

\begin{definition}[Box-stable exact recourse]\label{def:rec:qp-oracle}
Let $F_0(x)=\frac12x^{\mathsf T}Ax+a^{\mathsf T}x+c_0$ be a rational quadratic
on $[0,1]^n$ with core $C$ and residual $R$. A \emph{box-stable exact
recourse oracle} receives a rational core point $v\in[0,1]^k$, a rational
box $\prod_{i\in R}[\ell'_i,u'_i]\subseteq[0,1]^{n_R}$, possibly with
degenerate intervals, and rational linear coefficients $\gamma$, and returns
an exact rational global minimizer of $z\mapsto F_\gamma(v,z)$ on that box and
its value. Its bit work is bounded by a polynomial in the query length whose
exponent does not depend on $k$.
\end{definition}

An oracle that solves only the unrestricted residual problem does not
satisfy this definition, and the slab certificate cannot be run with it.

\begin{theorem}[Exact quadratic recourse]\label{thm:rec:qp}
Let $F_0(x)=\frac12x^{\mathsf T}Ax+a^{\mathsf T}x+c_0$ be a rational quadratic
on $X=[0,1]^n$, let $C$ be a supplied core of size $k\ge1$, and let $L>0$
be rational with $A_{ii}\le L$ for $i\in C$ (verified). Assume a box-stable
exact recourse oracle. There is a power of two $M$, computed from the base
instance with $\log_2M\le\poly(I)$, such that for the ambient perturbation
model with marginal law $U_{\sigma,M}$ the following hold.
\begin{enumerate}[label=(\roman*)]
\item On every draw the algorithm returns a rational global minimizer of
$F_\gamma$ on $X$ and the optimal value, of bit length $\poly(I)$.
\item The expected bit work, including oracle calls and the computation of
$M$, is at most $8^kQ_{\rm ex}\poly(I)$.
\end{enumerate}
No growth, uniqueness, nondegeneracy, or residual convexity premise is used.
\end{theorem}

The polynomial factor depends on the oracle's bound but has an exponent
independent of $k$. The algorithm runs the core search in exact mode; at
each level it tests \cref{lem:rec:exclusion} with at most $2n_R$ additional
oracle calls and then the tests of \cref{lem:rec:closure}. If they have not
passed by a base-computed level $J$, it enumerates the $3^n$ faces of the
box on the same draw (\cref{lem:rec:qp-faces}). The proof in
\cref{app:rec:main} shows that under point growth $g_0$ and active-gradient
margin $\tau$, both chosen from base data, all tests pass by level $J$, and
that the growth and active-gradient tails make the fallback rare. The
stopping depth depends only on base data: no rational-reconstruction bound
that grows with the sampling precision is needed.

\begin{corollary}[Residual forests]\label{cor:rec:forest}
In \cref{thm:rec:qp}, suppose that deleting the core $C$ from the
interaction graph of $A$ (vertices $[n]$, an edge $ij$ when $i\ne j$ and
$A_{ij}\ne0$) leaves a forest. Then the box-stable exact recourse oracle
exists, and the expected bit work is $8^kQ_{\rm ex}\poly(I)$. This bound is
fixed-parameter in the size $k$ of the supplied feedback vertex set and in
$L/\sigma$.
\end{corollary}

The oracle is the exact algorithm of \citet{DelPiaKhajavirad2026Forest} for
box-constrained quadratic programs whose interaction graph is a forest:
fixing core coordinates and adding linear terms change only linear
coefficients, and restricting a residual box and rescaling its nondegenerate
intervals keeps the residual graph a forest. The residual quadratic need not
be convex, have bounded negative inertia, or be separable. The corollary
is not fixed-parameter in treewidth alone: a graph of treewidth two can
require a large feedback vertex set, and box-constrained quadratic
programming is strongly NP-hard at treewidth two
\citep{DelPiaKhajavirad2026Forest}.

\subsection{Certified convex recourse for polynomial objectives}\label{sec:rec:poly}

For nonlinear recourse the conditional values are usually irrational, and
exact comparisons between them are unavailable. The theorem below uses only
certified intervals for conditional values and feasible rational
completions.

\begin{definition}[Box-stable certified recourse]\label{def:rec:cert-oracle}
For an explicit rational polynomial $F_0$ on $[0,1]^n$ with core $C$, a
\emph{box-stable certified recourse oracle} receives a rational core point
$v$, a rational residual box $Y'\subseteq[0,1]^{n_R}$, rational linear
coefficients $\gamma$, and a rational $\eta>0$. It returns a rational $\ell$
and a rational $z\in Y'$ with $\ell\le V_{Y'}(v)\le F_\gamma(v,z)\le\ell+\eta$.
Its bit work is polynomial in the query length, including the length of
$\eta$, with an exponent independent of $k$.
\end{definition}

\begin{lemma}[Certified convex recourse]\label{lem:rec:convex-oracle}
Fix $d$. If $\nabla^2_{RR}F_0\succeq0$ on $[0,1]^n$, a box-stable certified
recourse oracle exists with bit work $\poly_d$ in the query length. Moreover,
the returned lower bound can be taken as the tangent bound
\begin{equation}\label{eq:rec:tangent}
 \ell=f(z)+\min_{y\in Y'}\nabla f(z)^{\mathsf T}(y-z),\qquad
 f=F_\gamma(v,\cdot),
\end{equation}
which a verifier checks by evaluating $f$ and $\nabla f$ at the rational
point $z$ and choosing one endpoint per coordinate.
\end{lemma}

The proof (\cref{app:rec:convex}) uses weak convex optimization over a capped
epigraph \citep[Definition~2.1.10 and Corollary~4.2.7]{GroetschelLovaszSchrijver1988}
to find a feasible point of small objective error; convexity alone makes
\eqref{eq:rec:tangent} a global lower bound, and a smoothness bound along a
feasible segment makes its gap small. No strong convexity modulus is used.
The convexity premise is not decidable from an arbitrary polynomial in
general; it is certified, for example by an explicit sum of squares or by a
structural form such as a sum of even powers of affine functions.

\begin{theorem}[Polynomial recourse with certified conditional values]\label{thm:rec:poly}
Fix $d$. Let $F_0$ be an explicit rational polynomial of degree at most $d$
on $X=[0,1]^n$, let $C$ be a supplied core of size $k\ge0$, and let $L>0$ be
rational with \eqref{eq:rec:L} (verified by monomial bounds or certified).
Assume a box-stable certified recourse oracle; by
\cref{lem:rec:convex-oracle} one exists when $\nabla^2_{RR}F_0\succeq0$ on
$X$ (certified). There is a power of two $M$, computed from the base
instance with $\log_2M\le\poly_d(I)$, such that for the ambient perturbation
model with marginal law $U_{\sigma,M}$ the following hold.
\begin{enumerate}[label=(\roman*)]
\item On every draw the algorithm returns an exact global optimizer of
$F_\gamma$ on $X$. On ordinary draws the output is an implicit patch output
$(S,z_S,Q,g_0)$ of length $\poly_d(I)$, with a global proof record. On the
remaining draws it is the algebraic output of the fallback of
\cref{thm:count:fallback} on the same draw.
\item The expected bit work and the expected size of the global proof
record are at most $8^kQ_{\rm ap}\poly_d(I)$.
\item A $q$-bit evaluation of an implicit patch output costs
$\poly_d(I+q)$ bit operations; the expected cost of a $q$-bit evaluation of
the output, over all draws, is $\poly_d(I+q)$.
\end{enumerate}
No growth, uniqueness, or nondegeneracy premise is used.
\end{theorem}

The ordinary branch decides no exact comparison between irrational
conditional values. The patch descriptor determines one point because
$F_\gamma$ is strongly convex on $Q$, and the global proof record shows that
this point is the global optimizer; the record contains the pruning trace and
has only an expected size bound. An expanded algebraic representation of the
optimizer can have large degree even on an ordinary draw; the output contract
does not include one. For $k=0$ the problem is a convex polynomial on a box
with ambient noise. The theorem then still applies: the patch certificate
replaces any need to decide uniqueness or the active set.

The residual objective may be nonlinear, nonseparable, dense, and not
strongly convex. For example,
$F_0(v,z)=\psi(v)+\sum_iz_i^4+\norm{Dz-Ev}^2$ has convex recourse for any
rational matrices $D,E$ and any polynomial $\psi$. The core curvature bound
need not reflect the size of the core--residual coupling, and a small core
need not arise from a low-rank quadratic convexifier.

\begin{proposition}[A core not given by a low-rank convexifier]\label{prop:rec:rank}
For integers $m\ge1$ and rationals $\lambda,T>0$, $t_0\in[0,1]$,
$\beta\in\Q^m$, let
\[
 F(t,z)=\frac\lambda2(t-t_0)^2+T\,t\,P(z)-\beta^{\mathsf T}z,\qquad
 P(z)=\sum_{i=1}^mz_i^4+\Bigl(\sum_{i=1}^mz_i\Bigr)^4,
\]
on $[0,1]\times[0,1]^m$. Then $\partial_{tt}F=\lambda$, the residual problem
is convex for every $t$, and the core--residual cross derivatives grow
linearly with $T$. If a symmetric positive semidefinite matrix $K$ makes
$F+\frac12(t,z)^{\mathsf T}K(t,z)$ convex on the box, then $K_{zz}\succ0$; in
particular $\rank K\ge m$.
\end{proposition}

\Cref{thm:rec:poly} applies to this family with one core coordinate and
$L=\lambda$, so its numerical factor does not depend on $T$; the
coefficient $T$ enters only the polynomial factor through its bit length.
By the proposition, no fixed positive semidefinite quadratic correction of
rank independent of $m$ makes the objective convex. The proposition
excludes only fixed quadratic corrections. A one-coordinate correction
$Ct\log t$ with $C$ of order $T\max P$ does convexify the family, so the
example separates the theorem from small-rank quadratic convexifiers, not
from all difference-of-convex representations. It is not a hardness
result.

\subsection{Core-only noise and strongly convex recourse}\label{sec:rec:core-only}

\Cref{thm:rec:poly} perturbs every coordinate. Residual noise is what
produces full point growth and nonzero active gradients at residual bounds
with high probability. Under core-only noise neither holds in general: the
residual optimal set can be a segment on every draw, and residual bound
multipliers can vanish identically. A residual modulus restores closure.

\begin{theorem}[Core-only noise with strongly convex recourse]\label{thm:rec:core-only}
Fix $d$. Let $F_0$, $X=[0,1]^n$, $C$, and $L$ be as in \cref{thm:rec:poly},
with $k\ge1$ and $n_R\ge1$, and let $\mu>0$ be rational with
\begin{equation}\label{eq:rec:mu}
 \nabla^2_{RR}F_0(x)\succeq\mu I\qquad(x\in X)
\end{equation}
(certified). There is a power of two $M$, computed from the base instance
with $\log_2M\le\poly_d(I)$, such that for the core-only perturbation model
with marginal law $U_{\sigma,M}$ (all residual coefficients unperturbed) the
conclusions (i)--(iii) of \cref{thm:rec:poly} hold. The ordinary output is
an implicit patch output in which some original core and residual bounds may
be fixed.
\end{theorem}

No residual interiority, strict complementarity, active-set, or growth
premise is used, and residual bound multipliers may vanish identically. The
numerical factor contains the core curvature $L$ only. The modulus $\mu$ and
the size of mixed derivatives enter the polynomial factor through the bit
lengths of the derived constants. By \cref{prop:model:regret}, the exact
sampled optimizer is within $k\sigma$ of $\min_XF_0$ in original objective
value, independently of the number of residual variables.

The class admits a fixed rank-$k$ quadratic convexifier: the Schur complement
of the residual block and \eqref{eq:rec:K} give
$\nabla^2F_0+\alpha\diag(I_C,0)\succeq0$ for $\alpha=K_2+K_2^2/\mu$. Its
curvature can be much larger than $L$. The contribution of
\cref{thm:rec:core-only} is that the numerical factor keeps $L/\sigma$ under
core-only noise, despite changing and weakly active residual bounds; it is
not that the class lacks a low-rank convexifier.

The proof (\cref{app:rec:core-only}) has three new ingredients. First,
the residual selector $s(v)=\argmin_{z\in Y}F_0(v,z)$ is unique and
$(K_2/\mu)$-Lipschitz, the conditional value is continuously differentiable,
and projected growth $g_V$ of $V$ at its optimal core lifts to full point
growth $\min\{g_V/(1+2K_2^2/\mu^2),\mu/4\}$. The growth tail of
\cref{thm:count:growth-tail} applies to $V$ on the $k$-dimensional core box.
Second, core noise keeps the optimal core away from changes of the residual
active set, with a coefficient-independent probability bound. Third, small
residual multipliers can be released without losing a positive Hessian
modulus. We state the last two ingredients because they are the reason the
theorem needs no strict complementarity.

For a face $K$ of $[0,1]^k$ with $q\ge1$ free coordinates, identified with a
$q$-dimensional box, let $A_i^{\ell}=\{v\in K:s_i(v)=\ell_i\}$ and
$A_i^u=\{v\in K:s_i(v)=u_i\}$, and define
\begin{equation}\label{eq:rec:SK}
 S_K=\partial K\cup\bigcup_{i\in R}\bigl(\operatorname{bd}_KA_i^\ell\cup
 \operatorname{bd}_KA_i^u\bigr),
 \qquad
 \mathcal E_K=-\nabla_KV_0(S_K)\subseteq\R^q,
\end{equation}
where $\operatorname{bd}_K$ is the boundary relative to $K$, $V_0$ is the
conditional value without noise, and $\nabla_K$ is the gradient in the free
coordinates of $K$.

\begin{proposition}[Active-set changes are rarely near the optimum]\label{prop:rec:tube}
In the setting of \cref{thm:rec:core-only}, let $K$ be a face of $[0,1]^k$
with $q\ge1$ free coordinates.
\begin{enumerate}[label=(\alph*)]
\item $S_K$ and $\mathcal E_K$ are compact semialgebraic sets of dimension at
most $q-1$. If $a\in\operatorname{relint}K$ and
$0<\eta'<\dist(a,S_K)$, the relative ball of radius $\eta'$ about $a$ lies in
$\operatorname{relint}K$, and the bound status of each residual coordinate
of $s(v)$ is constant on it.
\item There is a base-computable $D_*$ with $\log D_*\le\poly_d(I)$ such that
$\mathcal E_K$ lies in the zero set of a nonzero polynomial of degree at most
$D_*$.
\item If $\gamma_K$ has independent coordinates with law $U_{\sigma,M}$, then
for every $\delta\ge0$
\[
 \Pr\{\dist(\gamma_K,\mathcal E_K)\le\delta\}
 \le\min\Bigl\{1,\ C(q,D_*)\Bigl(\frac\delta\sigma+\frac{\sqrt q}M\Bigr)\Bigr\},
 \qquad C(q,D)=16q^{q+1}D(4D+2)^{q-1}.
\]
\end{enumerate}
\end{proposition}

At an optimal core $a\in\operatorname{relint}K$ the free coordinates of the
noise satisfy $\gamma_K=-\nabla_KV_0(a)$, so (c) bounds the probability that
the optimal core is close to a change of the residual active set. Part~(b)
uses fixed-block quantifier elimination only as an analysis device: the
algorithm never computes $S_K$ or the polynomial. Part~(c) combines a tube
estimate for singular algebraic sets \citep{BasuLerario2023} with a jitter
argument that covers every atom of the grid, including atoms on
$\mathcal E_K$. The conclusion is a stable residual active set near the
optimum; it gives no positive lower bound on active multipliers.

\begin{proposition}[Releasing small residual multipliers]\label{prop:rec:release}
In the setting of \cref{thm:rec:core-only}, let $x^*=(a,s(a))$ with $a$ in
the relative interior of a face $K$ with $q\ge1$ free coordinates. Suppose the
residual bound status is constant on the relative ball of radius $\eta>0$
about $a$, $\nabla_KV_\gamma(a)=0$, and $\nabla_K^2V_\gamma(a)\succeq2gI$
with $g>0$. For an active residual coordinate $i$ let
$\lambda_i=\pm\partial_{z_i}F_0(x^*)\ge0$ be its inward multiplier. Put
$H=K_2/\mu$, $K_\lambda=\max\{1,K_3(1+H)^3\}$, and
\[
 \theta=\min\Bigl\{\frac{K_\lambda\eta^2}2,\ \frac{\mu g}{4n_RK_\lambda}\Bigr\},
 \qquad
 \nu=\min\Bigl\{\frac\mu2,\ \frac{\min\{g,\mu/2\}}{1+2H^2}\Bigr\}.
\]
Let $J$ be any set of active residual coordinates with $\lambda_i\le\theta$.
Then the Hessian of $F_0$ at $x^*$, restricted to the free coordinates of
$K$, the residual coordinates strictly inside their bounds, and $J$, is at
least $\nu I$. The same holds for every principal submatrix.
\end{proposition}

The bound uses the two-sided ball in the free core directions: nonnegativity
of a multiplier on both sides of $a$ and a Hessian bound force a small
multiplier to have a small gradient, which controls the cross terms between
released coordinates and the core. The algorithm never identifies the
active set or classifies multipliers. It fixes the coordinates whose
derivative signs are certified, which include all active residual
coordinates with $\lambda_i>\theta$ once the patch is small, and tests the
Hessian on the rest.

\begin{remark}[Interior recourse]\label{rem:rec:interior}
If, in addition, every conditional residual minimizer is interior to the
residual box (certified, for example by strict inward derivative signs on the
residual facets), then no residual coordinate is active. After the active
core coordinates are fixed, point growth alone gives a positive Hessian on
all remaining coordinates, and the proof of \cref{thm:rec:core-only}
needs neither \cref{prop:rec:tube} nor \cref{prop:rec:release}. No
quantitative distance from the residual bounds is used, and an accepted
patch may touch them.
\end{remark}

The following examples show which hypotheses cannot be dropped. Proofs are
in \cref{app:rec:examples}.

\begin{example}[Flat residual fibers]\label{ex:rec:fiber}
On $[0,1]^3$ with core $v$ and residual $(z_1,z_2)$ let
$F_\gamma(v,z)=(v-\frac12)^2+(z_1-vz_2)^2+\gamma v$. The residual Hessian is
positive semidefinite, $\partial_{vv}F_0\le4$, and for every
$\gamma\in(-1,1)$ the optimal core $a_\gamma=\frac12-\frac\gamma2$ is unique
with projected growth constant one. The optimal set
$\{(a_\gamma,a_\gamma s,s):s\in[0,1]\}$ is a segment, and the full Hessian is
indefinite at every point with $z_1\ne vz_2$. Hence no full-dimensional box
in $X$ carries a convex restriction of $F_\gamma$. Residual convexity without
a modulus does not support the closure of \cref{thm:rec:core-only}. If the
residual term is replaced by $(z_1-\frac12-v(z_2-\frac12))^2+(z_2-\frac12)^4$,
the residual problem is strictly convex with the unique interior minimizer
$(\frac12,\frac12)$ for every $v$, the optimizer
$(a_\gamma,\frac12,\frac12)$ is unique, and still the full Hessian is
indefinite at points arbitrarily close to it. Strict residual convexity and
interior uniqueness therefore do not replace the uniform modulus $\mu$.
\end{example}

\begin{example}[Vanishing multipliers]\label{ex:rec:weak}
On $[0,1]^2$ let $F_0(v,y)=\phi(v)+y^2$, with residual $y$. The residual
modulus is $2$, the selector is $s\equiv0$, and the multiplier of the active
bound $y=0$ vanishes identically, on every draw of core-only noise. The
coupled variant $F_0=(v-\frac12)^2+(y+(v-\frac12)^2)^2$ has multiplier
$2(v-\frac12)^2$, which vanishes at an interior core point while the
conditional value has quadratic growth. A closure that requires residual
multipliers bounded away from zero fails here; \cref{prop:rec:release} keeps
such coordinates free.
\end{example}

\begin{example}[Why active core coordinates are fixed first]\label{ex:rec:two-sided}
On $[0,1]^2$ let $F_\gamma(v,y)=\frac14v^2+vy+\frac12y^2+\gamma v$ with
residual $y$. The selector is $y=0$, $V_\gamma(v)=\frac14v^2+\gamma v$, the
residual modulus is one, and the full Hessian
$\bigl(\begin{smallmatrix}1/2&1\\1&1\end{smallmatrix}\bigr)$ is indefinite.
The multiplier $\lambda(v)=v$ is nonnegative on the feasible side only; at
the core vertex $v=0$ it vanishes while its derivative is one. For
$\gamma>0$ the optimum is $(0,0)$, the core derivative there is $\gamma$, and
the closure succeeds only because the sign test fixes $v=0$ first.
\Cref{prop:rec:release} applies only in the free directions of the optimal
core face.
\end{example}

\subsection{Scope}\label{sec:rec:scope}

All results in this section optimize the sampled objective exactly; none
recovers an optimizer of $F_0$. The core, the curvature bound $L$, the
residual convexity or modulus, and the box-stable oracle are inputs. The
results do not recognize a good core, and they do not decide residual
convexity of an arbitrary polynomial. The residual set must not depend on
the core (\cref{ex:rec:feasibility}), and the oracle must handle every
residual subbox. The numerical factors depend on the magnitude of
$L/\sigma$, so these are expected polynomial bounds only when that ratio
is polynomially bounded. The certificates are exact and sound on every draw,
while their early success, the stopping depth, and the size of the sampling
grid rest on growth and active-gradient tails for the base-selected finite
law (\cref{thm:count:growth-tail,lem:count:finite-tails}). Under core-only noise,
residual convexity without a modulus is not enough for the closure used here
(\cref{ex:rec:fiber}); it remains open whether a different certificate,
for instance one selecting an optimal residual face, gives a comparable
theorem. For integer residual variables the patch and Hessian tests are
replaced by label certificates; \cref{sec:int} treats that case, together
with core-only noise for integer flows.
